% Auto-generated / paper-ready snippet for orthogonal efficiency (C5 highlight).
% Cases and DR from orthogonal_reduction.csv; no timing columns in the table.

\subsection{Orthogonal Testing Efficiency}

\begin{figure*}[t]
\centering
% \includegraphics[width=0.85\linewidth]{figures/evaluation/orthogonal-reduction.png}
\caption{Orthogonal vs.\ exhaustive test-case counts as configuration items grow
under increasing service scales.}
\label{fig:orthogonal-reduction}
\end{figure*}

This experiment evaluates the efficacy of orthogonal test-case generation in
reducing the number of dynamic validation cases, and verifies that fault-detection
coverage is preserved after case merging.

We compare the orthogonal strategy with exhaustive validation, where each effect
is validated independently, across service scales
$N\in\{5,10,15,20,30,40,50\}$ and increasing configuration complexity
($C1$--$C5$).
We report the number of generated test cases and the reduction ratio.
To assess coverage preservation, we further compare fault detection rates on the
same fault set before and after orthogonal merging.

Figure~\ref{fig:orthogonal-reduction} shows how the number of generated test
cases grows with configuration items.
Exhaustive validation scales rapidly because each additional policy effect
introduces a separate validation path.
In contrast, the orthogonal strategy packs semantically independent, jointly
observable effects into shared executions and removes redundant cases.
As summarized in Table~\ref{tab:ortho-ablation}, we highlight settings with
high case reduction across different service scales and configuration sizes.
For example, with 20~services and 280 configuration items, the case count
drops from 100 to 25 (reduction $75.00\%$);
with 50~services and 1050 configuration items, it drops from 600 to 151
(reduction $74.83\%$).
Across these highlighted settings, the reduction stays in the $73\%$--$75\%$
range, while the detection rate remains $100\%$ both before and after merging,
because only independently observable effects are packed into a shared execution.
Thus orthogonal generation shrinks the validation suite without coverage loss.

\begin{table}[t]
\centering
\caption{High case-reduction settings across service scales and configuration sizes
(detection rate remains $100\%$ before and after merging).}
\label{tab:ortho-ablation}
\scriptsize
\begin{tabular}{ccccc}
\toprule
Services & cfg & Coverage (before) & Coverage (after) & Case reduction \\
\midrule
5 & 105 & 100\% & 100\% & 73.33\% \\
10 & 140 & 100\% & 100\% & 74.00\% \\
15 & 210 & 100\% & 100\% & 74.67\% \\
20 & 280 & 100\% & 100\% & 75.00\% \\
30 & 630 & 100\% & 100\% & 74.72\% \\
40 & 560 & 100\% & 100\% & 75.00\% \\
50 & 1050 & 100\% & 100\% & 74.83\% \\
\bottomrule
\end{tabular}
\end{table}
